\documentclass[11pt]{article}

\usepackage[margin=1in,top=1in,bottom=1in]{geometry}
\usepackage[utf8]{inputenc}
\usepackage[T1]{fontenc}
\usepackage{microtype}
\usepackage{float}
\usepackage{graphicx}
\usepackage{booktabs}
\usepackage{longtable}
\usepackage{amsmath}
\usepackage{amssymb}
\usepackage{array}
\usepackage{xcolor}
\usepackage{url}
\usepackage[colorlinks=true,linkcolor=blue,citecolor=blue,urlcolor=blue]{hyperref}
\usepackage{titlesec}
\usepackage{parskip}
\usepackage{caption}

% ── Figure and table numbering ─────────────────────────────────────────────
\renewcommand{\thefigure}{S\arabic{figure}}
\renewcommand{\thetable}{S\arabic{table}}
\renewcommand{\theequation}{S\arabic{equation}}

% ── Section formatting ─────────────────────────────────────────────────────
\titleformat{\section}{\large\bfseries}{SM\thesection.}{0.5em}{}
\titleformat{\subsection}{\normalsize\bfseries}{\thesubsection}{0.5em}{}

% ── Utility ────────────────────────────────────────────────────────────────
\newcommand{\pval}[1]{$p = #1$}
\newcommand{\betaval}[1]{$\beta = #1$}

\captionsetup[figure]{labelfont=bf,labelsep=period,font=small}
\captionsetup[table]{labelfont=bf,labelsep=period,font=small}

\title{\textbf{Supplementary Information}\\[6pt]
\large Genome streamlining is pervasive across prokaryotic functional categories,
but a distinctive cofactor-versus-resistance split links metal-gene investment
to ecological specialisation}

\author{Heather R.\ MacGregor$^{1,2}$ and Adam P.\ Arkin$^{2,3}$\\[4pt]
\small $^1$Department of Chemistry and Chemical Biology, University of California,
Berkeley, CA~94720, USA\\
\small $^2$Environmental Genomics and Systems Biology Division,
Lawrence Berkeley National Laboratory, Berkeley, CA~94720, USA\\
\small $^3$Department of Bioengineering, University of California,
Berkeley, CA~94720, USA}

\date{}

\begin{document}
\maketitle
\tableofcontents

% ============================================================================
\clearpage
\section{Gene list construction and evidence tiers}
\label{sm:genelist}
% ============================================================================

\subsection{Candidate KO universe}

The metal-gene list was assembled from three independent sources:

\begin{enumerate}
\item \textbf{KEGG ORTHOLOGY / BRITE hierarchies.} KOs classified under BRITE
  modules for heavy-metal resistance, metal transport, and metalloenzymes were
  downloaded in 2024.  Only KOs with an explicit metal annotation in the KEGG
  reaction or definition field were retained.

\item \textbf{BacMet2 database} (Pal et al.\ 2014).  Experimentally validated
  metal resistance genes from BacMet v2.0.  KEGG IDs were assigned by mapping
  BacMet gene names to KEGG gene identifiers via the KEGG REST API
  (\texttt{/find/genes/\{name\}} endpoint).

\item \textbf{FitnessBrowser} (Price et al.\ 2018).  Genome-wide transposon
  fitness screens under metallic stress conditions.  KOs with significant
  fitness defects ($|\text{fitness score}| > 1.5$, $|t| > 5$) under $\geq 1$
  metal stressor were included.
\end{enumerate}

KOs appearing in $\geq 2$ sources were cross-validated; those from only one
source were retained as lower-confidence candidates with appropriate tier labels.
The final list covers 730 KOs across 17 metals.

\subsection{Evidence tiers}

Each of the 730 candidate KOs was assigned to one of five evidence tiers:

\begin{table}[H]
\centering
\caption{Evidence tier definitions.}
\label{tab:tiers}
\small
\begin{tabular}{llrl}
\toprule
Tier & Label & $n$ KOs & Criterion \\
\midrule
1 & Tier~1 & 32 & Multi-source validated: BacMet2 \textsc{and} FitnessBrowser
  \textsc{and} KEGG module \\
2 & Tier~2 & 108 & Clear KEGG module definition; one additional source
  (BacMet2 or Fitness) \\
2-Fitness & Tier~2-Fitness & 116 & FitnessBrowser-validated only; no KEGG
  module annotation \\
3-BacMet & Tier~3-BacMet & 188 & BacMet2 literature-curated; no fitness
  evidence \\
3 & Tier~3 & 286 & KEGG BRITE only; ambiguous or multi-function annotation \\
\bottomrule
\end{tabular}
\end{table}

\textbf{Primary KO set} used in all confirmatory analyses: Tier~1 + Tier~2
= \textbf{140 KOs}.  These have the strongest multi-source evidence and
unambiguous KEGG functional annotations.  The full list with curation
rationale is provided in Supplementary Table~S2.

\subsection{Functional categories}

Each KO was assigned a primary functional category based on KEGG annotation
and BacMet2 gene function:

\begin{itemize}
\item \textbf{Transport / Homeostasis} — metal importers, exporters,
  metallochaperones, siderophore-uptake systems, and ferritins.

\item \textbf{Resistance / Detoxification} — metal-efflux pumps, detoxification
  enzymes, and metal-responsive regulatory proteins that confer resistance.

\item \textbf{Sensing / Regulation} — two-component sensor kinases, MerR-family
  metalloregulators, and other metal-sensing regulatory elements.

\item \textbf{Cofactor Biosynthesis} — enzymes in Fe--S cluster assembly,
  molybdopterin (Moco) biosynthesis, and cobalamin pathway steps that
  incorporate metal cofactors.

\item \textbf{Metal-dependent Metabolism} — metabolic enzymes whose catalytic
  activity requires a metal cofactor, supported by KEGG module assignment and
  fitness screen evidence.
\end{itemize}

Assignment is based on the most specific category for each KO;
\texttt{overlap\_flag = True} identifies six KOs with dual-function ambiguity
in the full 730-KO list (not in the primary 140-KO set).

\subsection{Metal coverage}

The primary 140-KO set covers 17 metals: Fe, Zn, Co, Mn, Cu, Ni, S, Mo, Cd,
Al, Tl, Ag, Hg, Au, Pb, As, and Bi.  Metal annotations are listed per KO in
Supplementary Table~S2.  Metal-specific PGLS subsets were defined by filtering
to KOs whose \texttt{metals} field contains the target metal.  Nine metals with
$\geq 34$ primary KOs were included in the metal-specific confirmatory analysis:
Co (101 KOs), Fe (98), Ni (84), S (74), Cu (71), Zn (67), Tl (45), Al (34),
Mn (34). For Al, Tl, and S --- elements with fewer dedicated resistance systems --- KOs were assigned based on fitness-screen evidence or documented biochemical function; full curation rationale for these three metals is provided in Supplementary Table~S4.

\subsection{Pfam / InterPro structural QC}

To assess whether the primary KO set has independent structural evidence of
metal interaction, we performed a three-step API pipeline for all 140 KOs:

\begin{enumerate}
\item \textbf{KO $\to$ gene:} KEGG REST \texttt{link/genes/ko:\{KO\}} filtered
  to representative model organisms in priority order (\textit{E.\ coli},
  \textit{B.\ subtilis}, \textit{P.\ aeruginosa}, \textit{Synechocystis},
  \textit{M.\ tuberculosis}, \textit{S.\ coelicolor}, \textit{R.\ prowazekii},
  \textit{C.\ jejuni}, \textit{B.\ pertussis}, \textit{M.\ genitalium});
  then \texttt{conv/uniprot/\{gene\_id\}} to obtain a UniProt accession.
  Rate: 0.4\,s per call; 3-retry exponential back-off.

\item \textbf{Gene $\to$ Pfam:} UniProt REST
  \texttt{/\{acc\}?fields=xref\_pfam\&format=tsv} to retrieve Pfam domain IDs.

\item \textbf{Pfam $\to$ clan:} InterPro API
  \texttt{/entry/pfam/\{PF\}/?format=json} $\to$
  \texttt{metadata.set\_info.accession} to identify the InterPro clan.
  Rate: 0.3\,s per call.
\end{enumerate}

Metal-binding domains were defined by five InterPro clans (verified in
InterPro version 97+): CL0704 (HMA heavy-metal-associated domain), CL0344
(4Fe--4S ferredoxin), CL0486 (Fer2 / 2Fe--2S ferredoxin), CL0361 (C2H2 zinc
finger), and CL0193 (MBB metal beta-barrel); plus $\approx$25 metal-binding
singleton Pfams identified from literature.

\textbf{Results:} 10/140 (7.1\%) KOs have $\geq 1$ metal-binding clan or
singleton Pfam.  The 130 without clan evidence fall into three classes:
\texttt{pfam\_no\_metal\_clan} (113 KOs — Pfam domains found but none in
metal-binding clans; predominantly ABC transporter ATPase / permease scaffolds
PF00005 / PF00950 / PF01032, MerR-family sensors, and outer-membrane efflux
proteins); \texttt{no\_gene\_found} (15 KOs — no representative bacterial gene
in KEGG REST); and \texttt{no\_pfam\_domains} (2 KOs: ARN / K08197 and cbiL /
K16915).  The low InterPro coverage is expected: metal transporters act via
substrate selectivity at the binding pocket rather than via catalytic
metal-cofactor domains, and will systematically fail clan-based Pfam filtering
(see also Fig.~S14 and §5.1 of the main text).

% ============================================================================
\section{Per-megabase KO density calculation}
\label{sm:density}
% ============================================================================

\subsection{Genomic data sources}

\textbf{Primary dataset (bacteria, $n = 1{,}574$ genera):} BERDL Spark cluster,
namespace \texttt{kbase.ke\_pangenome}.  The \texttt{genus\_clusters} table
contains one row per GTDB genus with pangenome statistics (total gene count,
mean genome size in bp, and the full KO annotation set across all genomes in
the genus-level cluster).  Phylum / kingdom annotations joined from
\texttt{kbase.ke\_pangenome.gtdb\_taxonomy\_r214v1}.

\textbf{Geochemical replication dataset (soil concentration,
$n = 482$ genera):} \texttt{refdata.ngsa} (Australian National Geochemical
Survey of Australia; Geoscience Australia).  KO content extracted by the same
pipeline, restricted to genera detected in Australian soil and sediment samples.
Soil metal concentrations (ICP-MS, mg\,kg$^{-1}$) joined by a 200\,km radius
spatial join between NGSA sampling sites and genus occurrences.

\subsection{KO density formula}

For each genus $g$:

\begin{equation}
\text{ko\_per\_mb\_primary}(g) =
  \frac{n_{\text{distinct primary KOs}}(g)}{\text{mean\_genome\_size\_Mb}(g)}
\end{equation}

where $\text{mean\_genome\_size\_Mb} = \text{mean\_genome\_size\_bp} / 10^6$
averaged across all genomes in the genus cluster, and
$n_{\text{distinct primary KOs}}$ counts unique KO IDs from the 140-KO primary
set present in $\geq 1$ genome of the genus.

\subsection{Predictor standardisation}

For PGLS, \texttt{ko\_per\_mb\_primary} is $z$-scored within each analysis
dataset:
\begin{equation}
\text{predictor\_z} = \frac{\text{ko\_per\_mb\_primary} -
  \mu(\text{ko\_per\_mb\_primary})}{\sigma(\text{ko\_per\_mb\_primary})}
\end{equation}

Standardisation is performed separately within each analysis dataset
(bacteria, archaea, geochemical replication), so $\beta$ values are not
directly comparable across datasets.

\subsection{Primary dataset summary}

\begin{table}[H]
\centering
\caption{Summary statistics for the primary bacterial dataset ($n = 1{,}574$
genera).}
\label{tab:datasummary}
\small
\begin{tabular}{ll}
\toprule
Statistic & Value \\
\midrule
$n$ genera & 1,574 \\
\texttt{ko\_per\_mb\_primary}: mean (SD) & 8.61 (3.83) KOs/Mb \\
\texttt{ko\_per\_mb\_primary}: range & 0.78--34.56 KOs/Mb \\
Mean genome size: mean (SD) & 3.73 (1.63) Mb \\
Dominant phyla & Proteobacteria (43\%), Firmicutes (21\%), \\
 & Actinobacteria (13\%), Bacteroidetes (12\%) \\
\bottomrule
\end{tabular}
\end{table}

% ============================================================================
\section{Ecological niche breadth quantification}
\label{sm:niche}
% ============================================================================

\subsection{Levins' standardised niche breadth}

Niche breadth was quantified using Levins' $B$ (Levins 1968), standardised
following Colwell \& Futuyma (1971):

\begin{equation}
B = \frac{1}{\sum_i p_i^2} \qquad\quad
B_\mathrm{std} = \frac{B - 1}{n - 1}
\end{equation}

where $p_i$ is the proportional abundance of a genus in sample $i$ (relative
to total abundance of that genus across all samples) and $n$ is the total
number of samples.  $B_\mathrm{std} \in [0,1]$: 0 = complete specialist
(all occurrences in one sample), 1 = complete generalist (perfectly even
distribution).  Implementation: \texttt{scripts/niche\_utils.py::levins\_b\_std()}.

\subsection{Primary niche breadth dataset (MicrobeAtlas)}

\textbf{Data source:} \texttt{arkinlab.microbeatlas} (BERDL Spark) — genus-level
OTU abundance matrix across $>400{,}000$ global 16S rDNA amplicon samples
aggregated in the MicrobeAtlas database (Yilmaz et al.).  Each OTU linked to a
genus via \texttt{otu\_pangenome\_link\_v2.csv}; per-genus $B_\mathrm{std}$
computed as the mean across all OTUs assigned to that genus.

\textbf{Minimum sample filter:} Genera detected in $<5$ samples are excluded.
Sensitivity checks at $n_\mathrm{min} = 50$ and $n_\mathrm{min} = 150$
produce identical results (analyses S4 and S5, Table S2 of the main text),
confirming that most genera in the dataset exceed these thresholds.

\subsection{EMP niche breadth (secondary, exploratory)}

\textbf{Data source:} Earth Microbiome Project 16S data
(\texttt{refdata.emp\_16s}).  EMPO (Earth Microbiome Project Ontology)
level-3 habitat categories used as the niche axis; level-2 used as fallback
when $<5$ distinct EMPO-3 categories were present for a genus.  Levins
$B_\mathrm{std}$ computed from the genus $\times$ EMPO-category co-occurrence
matrix.  Minimum: $\geq 5$ EMP samples per genus.  $n = 539$ genera after
merge with primary KO density data.

\textbf{Result:} $\lambda = 0.055$ (little phylogenetic signal in EMP niche
breadth), $\beta = -0.019$, $\text{SE} = 0.0115$, $p = 0.099$ (non-significant,
directionally consistent; see Fig.~S13).

\subsection{Soil-specialist definition}

For the soil-restricted sensitivity analysis (analysis S7 of the main text),
genera were classified as soil specialists via MicrobeAtlas \texttt{Env\_Level\_1}
using ecosystem categories: \textit{soil}, \textit{agricultural}, \textit{farm},
\textit{field}, \textit{paddy}, \textit{peatland}, \textit{desert}, \textit{shrub}.
Procedure: (1) identify the dominant \texttt{Env\_Level\_1} category per OTU
(row with maximum \texttt{n\_samples\_detected}); (2) flag OTU as
\texttt{is\_soil} if dominant environment $\in$ SOIL\_ENVS; (3) join OTUs to
genera; (4) compute $\text{frac\_soil} = n_\text{soil OTUs} / n_\text{OTUs}$
per genus; (5) classify genus as soil specialist if $\text{frac\_soil} > 0.5$.
Soil-restricted dataset: $n = 162$ genera.

% ============================================================================
\section{Phylogenetic tree}
\label{sm:tree}
% ============================================================================

\textbf{Source:} GTDB (Genome Taxonomy Database) Release 214v1.  The full
genus-level phylogeny was downloaded in Newick format and pruned to retain only
genera present in the PGLS input dataset using
\texttt{dendropy.Tree.prune\_taxa\_with\_labels()}.

\begin{itemize}
\item Bacterial tree: \texttt{data/gtdb\_bac\_genus\_pruned.tree}
  (2,283 taxa before pruning; 1,574 retained for the primary analysis).

\item Archaeal tree: \texttt{data/gtdb\_arc\_genus\_pruned.tree}
  (pruned to 95 genera for the archaeal replication; 41 for the archaeal
  sensitivity analysis at $n_\text{min} = 5$ genera per taxon).
\end{itemize}

Tip labels use lowercase genus names with spaces replaced by underscores,
matching the \texttt{genus\_lower} column in all data files.  All branch
lengths are in units of substitutions per site.  The circular cladogram of
the bacterial tree coloured by KO density and niche breadth is shown in
Fig.~S1.

% ============================================================================
\section{Phylogenetic generalised least squares}
\label{sm:pgls}
% ============================================================================

\subsection{Implementation}

All PGLS models were fitted using \texttt{scripts/pgls\_utils.py}, a custom
implementation using \texttt{dendropy} for tree operations and
\texttt{scipy} / \texttt{numpy} for optimisation and linear algebra.

\textbf{Variance--covariance matrix.}
Built from the pruned GTDB tree as $V[i,j] = $ shared branch length from root
to MRCA$(i,j)$ (= root-to-tip path length for $i = j$).

\textbf{Pagel's $\lambda$.}
Off-diagonal elements of $V$ are scaled by $\lambda \in (0,1)$:
\begin{equation}
V_\lambda[i,j] = \lambda \cdot V[i,j] \quad (i \neq j); \qquad
V_\lambda[i,i] = V[i,i]
\end{equation}
$\lambda = 0$: OLS (no phylogenetic signal); $\lambda = 1$: Brownian motion
(full phylogenetic signal).  $\lambda$ is estimated by maximising the profile
log-likelihood:
\begin{equation}
\log L = -\tfrac{1}{2}\bigl[n\log(2\pi) + n\log(\sigma^2) +
  \log|V_\lambda| +
  (\mathbf{y} - \mathbf{X}\boldsymbol{\beta})^\top
  V_\lambda^{-1}
  (\mathbf{y} - \mathbf{X}\boldsymbol{\beta}) / \sigma^2\bigr]
\end{equation}
Optimisation via \texttt{scipy.optimize.minimize\_scalar()} with bounds
$(10^{-4},\, 1 - 10^{-4})$ and \texttt{method='bounded'}.

\textbf{GLS fit.}
After $\lambda$ optimisation, GLS coefficients $\hat{\boldsymbol{\beta}}$ and
standard errors are estimated:
\begin{align}
\hat{\boldsymbol{\beta}} &=
  (\mathbf{X}^\top V^{-1} \mathbf{X})^{-1} \mathbf{X}^\top V^{-1} \mathbf{y}
\\[4pt]
\sigma^2 &=
  \frac{(\mathbf{y} - \mathbf{X}\hat{\boldsymbol{\beta}})^\top
    V_\lambda^{-1}
    (\mathbf{y} - \mathbf{X}\hat{\boldsymbol{\beta}})}{n - p}
\\[4pt]
\mathrm{SE}(\hat{\beta}_j) &=
  \sqrt{\bigl[(\mathbf{X}^\top V^{-1} \mathbf{X})^{-1}\bigr]_{jj} \cdot \sigma^2}
\end{align}
where $p$ is the number of predictors including the intercept.

\textbf{Model statistics:}
$t_j = \hat{\beta}_j / \mathrm{SE}(\hat{\beta}_j)$ with $\mathrm{df} = n - p$;
$\mathrm{AIC} = -2\log L + 2(p + 1)$;
$\Delta\mathrm{AIC} = \mathrm{AIC}(\text{model}) - \mathrm{AIC}(\text{null})$;
$r^2 = 1 - \mathrm{Var}(\text{residuals}) / \mathrm{Var}(\text{response})$.

\subsection{Primary result}

\begin{table}[H]
\centering
\caption{Primary PGLS result (bacteria, $n = 1{,}574$ genera; MicrobeAtlas
niche breadth vs Tier~1+2 KO density).}
\label{tab:p1}
\small
\begin{tabular}{ll}
\toprule
Parameter & Value \\
\midrule
$n$ genera (bacteria) & 1,574 \\
Pagel's $\lambda$ & 0.757 \\
$\beta$ (\texttt{ko\_per\_mb\_primary}, $z$-scored) & $-0.0207$ \\
SE & $0.00368$ \\
$t$-statistic & $-5.629$ \\
$p$ (raw) & $2.14 \times 10^{-8}$ \\
$p$ (FDR, joint with P2 and P3) & $4.28 \times 10^{-8}$ \\
$r^2$ & 0.046 \\
$\Delta$AIC vs null & $-29.41$ \\
\midrule
Archaeal replication ($n = 95$): $\beta$ & $-0.0137$ ($p = 0.119$, FDR $= 0.178$) \\
Geochemical replication ($n = 482$): $\beta$ & $-0.00171$ ($p = 0.755$) \\
\bottomrule
\end{tabular}
\end{table}

% ============================================================================
\section{Multiple testing correction}
\label{sm:mtc}
% ============================================================================

All corrections use the Benjamini--Hochberg (BH) procedure.

\begin{description}
\item[Primary three-way replication (§2.1 of main text).]
  P-values from the three primary confirmatory tests (bacteria, archaea,
  geochemical) corrected jointly ($n = 3$).
  FDR: bacteria $= 4.28 \times 10^{-8}$; archaea $= 0.178$;
  geochemical $= 0.755$.

\item[Evidence-tier comparisons (§2.2 of main text).]
  Two tier comparisons (\texttt{all\_non\_ambiguous} and \texttt{bacmet\_only})
  corrected together ($n = 2$).

\item[Functional category comparisons (§3 of main text).]
  Five functional category PGLS models corrected together ($n = 5$).

\item[Metal-specific analyses (§4 of main text).]
  Nine metal-specific PGLS models corrected together ($n = 9$);
  all nine pass FDR $< 0.05$.

\item[Clade-stratified analyses.]
  Four phylum-stratified PGLS models corrected together ($n = 4$).
\end{description}

No multiple testing correction was applied to pre-specified confounder checks
(§5.2 of the main text) or to the eight sensitivity analyses (§5.3), as these
are robustness tests of a single primary result rather than independent
hypotheses.

% ============================================================================
\section{Confounder analyses}
\label{sm:confounders}
% ============================================================================

For each potential confounder, we refitted the PGLS model adding the confounder
as a second predictor alongside \texttt{predictor\_z} and recorded the change
in the primary $\beta$:
\begin{equation}
\Delta\beta\,(\%) =
  \frac{|\hat{\beta}_{\text{with confounder}} - \hat{\beta}_{\text{baseline}}|}
       {|\hat{\beta}_{\text{baseline}}|}
  \times 100
\end{equation}
\textbf{Decision rule:} $\Delta\beta > 50\%$ \textsc{and} $p > 0.05$ for the
density predictor $\Rightarrow$ confounder explains the effect.

\begin{table}[H]
\centering
\caption{Confounder analysis results.}
\label{tab:confounders}
\small
\begin{tabular}{p{3.2cm}p{4.5cm}rp{4cm}}
\toprule
Confounder & Operationalisation & $\Delta\beta$ (\%) & Verdict \\
\midrule
Genome size & Mean genome size in Mb ($z$-scored) & 46.7 &
  ROBUST — below threshold; $p = 0.006$ \\
GC content & Mean genomic GC\% ($z$-scored) & 23.7 &
  ROBUST \\
Isolation source & Categorical dummy (culture environment) & 14.5 &
  ROBUST \\
Mean latitude & Mean absolute latitude per genus & 51.8 &
  ATTENUATED — $\beta$ becomes \emph{more} negative ($-0.031$), inconsistent
  with confounding; latitude amplifies rather than explains the signal \\
Dominant biome & Majority \texttt{Env\_Level\_1} (MicrobeAtlas) & 5.8 &
  ROBUST \\
\bottomrule
\end{tabular}
\end{table}

No confounder meets the pre-specified decision rule ($\Delta\beta > 50\%$
\textsc{and} non-significant).  Genome size partially attenuates $\beta$ by
46.7\% but the predictor remains significant ($p = 0.006$), suggesting that
genome size is a partial mediator rather than a confounder.

% ============================================================================
\section{Sensitivity analyses}
\label{sm:sensitivity}
% ============================================================================

Eight pre-specified sensitivity analyses tested robustness to modelling
assumptions.  All used the same response variable (\texttt{mean\_levins\_B\_std})
and the same primary KO density predictor, with the modification noted.

\begin{table}[H]
\centering
\caption{Sensitivity analysis results.  Six of eight analyses are
  directionally consistent with the primary result and statistically
  significant.}
\label{tab:sensitivity}
\small
\begin{tabular}{clrrl}
\toprule
ID & Modification & $\beta$ & $p$ & Notes \\
\midrule
S1 & $\lambda$ fixed $= 0$ (OLS) & $-0.0319$ & $\approx 0$ &
  Stronger effect under OLS \\
S2 & $\lambda$ fixed $= 1$ (Brownian) & $-0.0182$ & $3.66 \times 10^{-6}$ &
  Consistent \\
S3 & Archaea, $n_{\min} = 5$ genera & $-0.0183$ & $0.084$ &
  $n = 41$; directionally consistent, underpowered \\
S4 & $n_{\min} = 50$ MicrobeAtlas samples & $-0.0207$ & $2.14 \times 10^{-8}$ &
  Identical to primary \\
S5 & $n_{\min} = 150$ MicrobeAtlas samples & $-0.0207$ & $2.14 \times 10^{-8}$ &
  Identical to primary \\
S6 & Raw (unstandardised) Levins $B$ & $-0.286$ & $1.11 \times 10^{-11}$ &
  Stronger signal; same direction \\
S7 & Australia-only subset & $-0.00171$ & $0.755$ &
  Null; same result as geochemical replication \\
S8 & Northern hemisphere only & $-0.0302$ & $3.24 \times 10^{-6}$ &
  Stronger than primary; consistent \\
\bottomrule
\end{tabular}
\end{table}

S4 and S5 produce identical results because all 1,574 genera in the primary
dataset have MicrobeAtlas detection counts well above 150, making both
thresholds redundant.  S7 is numerically equivalent to the geochemical
replication (same $n = 482$ Australian subset; Australia is represented in
the global dataset but the narrowed geographic range removes the global
diversity signal).

% ============================================================================
\section{Pre-registration and analysis governance}
\label{sm:prereg}
% ============================================================================

\subsection{Pre-registration}

All confirmatory analyses (primary PGLS, archaeal replication, geochemical
replication, evidence-tier comparisons, functional category comparisons,
metal-specific comparisons, and all eight sensitivity analyses) were
pre-specified in \texttt{RESEARCH\_PLAN.md} prior to data analysis.  This
document recorded the hypothesis direction ($\beta < 0$), the response variable
(\texttt{mean\_levins\_B\_std}), the predictor (\texttt{ko\_per\_mb\_primary}),
the tree source (GTDB r214v1), the $\lambda$ estimation method (ML), the
multiple testing procedure (BH), and the decision rule for confounders
($\Delta\beta > 50\%$ \textsc{and} $p > 0.05$).

Exploratory analyses (EMP niche breadth validation, BacDive geographic breadth,
ENIGMA site-level replication, latitude mechanism tests, inverse PGLS, RDA
variance partitioning, and community-weighted mean prediction) are labelled as
such throughout the manuscript and were not pre-specified.

\subsection{Analysis governance}

Every reported statistic traces to a row in \texttt{INTERPRETATION\_TABLE.md},
which records the notebook cell ID, output CSV file, and direction prediction
for each confirmatory test.  Statistical values were entered from saved CSV
files in \texttt{data/}, not from notebook-cell display output.  Exploratory
analyses are tracked in a separate section of the same table, labelled
``(Exploratory)''.

% ============================================================================
\section{Latitude mechanism tests}
\label{sm:latitude}
% ============================================================================

The latitude amplification result (adding mean absolute latitude as a covariate
increases $|\beta|$ from $-0.021$ to $-0.031$, §5.2 of main text) motivated a
set of nine mechanistic tests examining whether a specific environmental
variable could explain the latitudinal gradient in niche breadth.  These are
exploratory analyses.

Models A--I each added a single environmental predictor alongside latitude to
the primary PGLS, testing whether the latitude effect was absorbed by:

\begin{itemize}
\item \textbf{Model A (bedrock Cr composite)}: GeoROC lithospheric
  metal-index principal component (driven by Cr: BH $p = 6.7 \times 10^{-9}$,
  and Co: BH $p = 0.0084$).

\item \textbf{Model B (temperature range)}: CHELSA annual temperature range
  (°C).

\item \textbf{Model C--H}: Ore-deposit proximity; bioavailable metal mobility
  index; soil metal hazard (CSU classification); soil temperature; lithology
  class; REE deposit proximity.

\item \textbf{Model I (soil organic matter)}: SOM (\%) as a redox proxy.
  \textbf{Key result:} SOM is an independent negative predictor
  ($\beta = -0.016$, $p = 6.8 \times 10^{-5}$); absorbs the pH association
  (pH becomes non-significant when SOM is included), consistent with
  organic-matter-driven metal speciation and bioavailability.

\item \textbf{Redox controls}: Soil moisture and SOM do not attenuate the
  Cr or Co bedrock signals ($p = 0.44$ and $0.50$ for attenuation tests),
  arguing against a waterlogging / soil-reduction mechanism for the serpentine
  Cr effect.
\end{itemize}

All 28 extended models are in \texttt{data/latitude\_mechanism\_results.csv}.

% ============================================================================
\clearpage
\section*{Supplementary Tables}
\addcontentsline{toc}{section}{Supplementary Tables}
\setcounter{table}{1}
% ============================================================================

% -- Table S2: 140-KO primary metal gene list ---------------------------------

\begin{longtable}{@{}llllp{3.8cm}l@{}}
\caption{Primary 140-KO metal-gene list with evidence tier, functional category,
metal targets, and source databases.  Tier: T1 = Tier~1 (multi-source validated,
$n = 32$); T2 = Tier~2 (KEGG module + one additional source, $n = 108$).
Category abbreviations: Transport = Transport/Homeostasis;
Resistance = Resistance/Detoxification; Sensing = Sensing/Regulation;
Cofactor = Cofactor Biosynthesis; Metabolism = Metal-dependent Metabolism.
Source: KEGG = KEGG module annotation; BacMet2 = BacMet v2.0 database;
Fitness = FitnessBrowser transposon screen.}
\label{tab:ko_list}\\
\toprule
KO & Gene & Tier & Category & Metals & Source \\
\midrule
\endfirsthead
\multicolumn{6}{l}{\small\textit{Table S2 continued.}}\\
\toprule
KO & Gene & Tier & Category & Metals & Source \\
\midrule
\endhead
\midrule
\multicolumn{6}{r}{\small\textit{Continued on next page.}}\\
\endfoot
\bottomrule
\endlastfoot
K01534 & zntA & T1 & Transport & Zn, Cd, Cu, Ni, Co, Hg, Pb & BacMet2, Fitness \\
K02006 & cbiO & T2 & Transport & Co, Ni & --- \\
K02007 & cbiM & T2 & Transport & Co, Ni & --- \\
K02008 & cbiQ & T2 & Transport & Co, Ni & --- \\
K02009 & cbiN & T2 & Transport & Co, Ni & --- \\
K02010 & afuC & T1 & Transport & Fe, Co, Zn & BacMet2, Fitness \\
K02012 & afuA & T1 & Transport & Fe, Co, Zn & BacMet2, Fitness \\
K02013 & ABC.FEV.A & T2 & Transport & Fe & --- \\
K02015 & ABC.FEV.P & T2 & Transport & Fe & --- \\
K02016 & ABC.FEV.S & T1 & Transport & Fe & BacMet2 \\
K02068 & STAR1 & T1 & Transport & Fe, Ni & BacMet2, Fitness \\
K02069 & STAR2 & T1 & Transport & Fe, Ni & BacMet2, Fitness \\
K02074 & ABC.ZM.A & T2 & Transport & Mn, Zn & --- \\
K02075 & ABC.ZM.P & T2 & Transport & Mn, Zn & --- \\
K02077 & ABC.ZM.S & T2 & Transport & Mn, Zn & --- \\
K02190 & cbiK & T2 & Transport & Co & KEGG \\
K02230 & cobN & T2 & Transport & Co & KEGG \\
K03523 & bioY & T2 & Transport & --- & --- \\
K03795 & cbiX & T2 & Transport & Co & KEGG \\
K05895 & cobK-cbiJ & T2 & Transport & Co & KEGG \\
K08197 & ARN & T2 & Transport & Fe & --- \\
K08225 & entS & T2 & Transport & Fe & --- \\
K08970 & rcnA & T1 & Transport & Co, Ni & BacMet2, Fitness \\
K09815 & znuA & T1 & Transport & Zn & BacMet2 \\
K09816 & znuB & T1 & Transport & Zn & BacMet2 \\
K09817 & znuC & T1 & Transport & Zn & BacMet2 \\
K09818 & ABC.MN.S & T2 & Transport & Mn, Fe & --- \\
K09819 & ABC.MN.P & T2 & Transport & Mn, Fe & --- \\
K09820 & ABC.MN.A & T2 & Transport & Mn, Fe & --- \\
K09882 & cobS & T2 & Transport & Co & KEGG \\
K09883 & cobT & T2 & Transport & Co & KEGG \\
K10094 & cbiK & T2 & Transport & Ni & --- \\
K10830 & psaB & T2 & Transport & Mn, Zn & --- \\
K10834 & hmuV & T2 & Transport & Fe & --- \\
K11601 & mntC & T2 & Transport & Mn & --- \\
K11602 & mntB & T2 & Transport & Mn & --- \\
K11603 & mntA & T2 & Transport & Mn & --- \\
K11604 & sitA & T1 & Transport & Mn, Fe & BacMet2 \\
K11605 & sitC & T1 & Transport & Mn, Fe & BacMet2 \\
K11606 & sitD & T1 & Transport & Mn, Fe & BacMet2 \\
K11607 & sitB & T1 & Transport & Mn, Fe & BacMet2 \\
K11704 & mtsA & T2 & Transport & Mn, Cu, Fe, Zn & --- \\
K11705 & mtsC & T2 & Transport & Mn, Cu, Fe, Zn & --- \\
K11706 & mtsB & T2 & Transport & Mn, Cu, Fe, Zn & --- \\
K11707 & troA & T1 & Transport & Mn, Fe, Zn & BacMet2 \\
K11708 & troC & T1 & Transport & Mn, Fe, Zn & BacMet2 \\
K11709 & troD & T1 & Transport & Mn, Fe, Zn & BacMet2 \\
K11710 & troB & T1 & Transport & Mn, Fe, Zn & BacMet2 \\
K12346 & SMF & T2 & Transport & Fe & --- \\
K14193 & isdA & T2 & Transport & Fe & --- \\
K14709 & SLC39A1\_2\_3 & T2 & Transport & Zn & --- \\
K14710 & SLC39A4 & T2 & Transport & Zn & --- \\
K14711 & SLC39A5 & T2 & Transport & Zn & --- \\
K14712 & SLC39A6 & T2 & Transport & Zn & --- \\
K14713 & SLC39A7 & T2 & Transport & Zn & --- \\
K14714 & SLC39A8 & T2 & Transport & Zn & --- \\
K14715 & SLC39A9 & T2 & Transport & Zn & --- \\
K14716 & SLC39A10 & T2 & Transport & Zn & --- \\
K14717 & SLC39A11 & T2 & Transport & Zn & --- \\
K14718 & SLC39A12 & T2 & Transport & Zn & --- \\
K14719 & SLC39A13 & T2 & Transport & Zn & --- \\
K14720 & SLC39A14 & T2 & Transport & Zn & --- \\
K16302 & CNNM & T2 & Transport & --- & --- \\
K16783 & bioN & T2 & Transport & --- & --- \\
K16784 & bioM & T2 & Transport & --- & --- \\
K16915 & cbiL & T2 & Transport & Ni & --- \\
K19971 & psaA & T2 & Transport & Mn, Zn & --- \\
K19972 & psaC & T2 & Transport & Mn, Zn & --- \\
K19973 & mntA & T2 & Transport & Mn & --- \\
K19975 & mntC & T2 & Transport & Mn & --- \\
K19976 & mntB & T2 & Transport & Mn & --- \\
K20170 & nicA1 & T2 & Transport & --- & KEGG \\
K20982 & LOVC & T2 & Transport & --- & KEGG \\
K22011 & cfbA & T2 & Transport & Co & KEGG \\
K23185 & fepB & T2 & Transport & --- & --- \\
K23186 & fepD & T2 & Transport & Fe & --- \\
K23187 & fepG & T2 & Transport & Fe & --- \\
K23188 & fepC & T2 & Transport & Fe & --- \\
K24175 & MFSD5 & T2 & Transport & Mo & --- \\
K25027 & btuC & T2 & Transport & Co & --- \\
K25028 & btuD & T2 & Transport & Co & --- \\
K25034 & btuF & T2 & Transport & Co & --- \\
K25109 & sirA & T2 & Transport & Fe & --- \\
K25110 & sirB & T2 & Transport & Fe & --- \\
K25111 & sirC & T2 & Transport & Fe & --- \\
K25113 & htsA & T2 & Transport & Fe & --- \\
K25114 & htsB & T2 & Transport & Fe & --- \\
K25115 & htsC & T2 & Transport & Fe & --- \\
K25116 & isdD & T2 & Transport & Fe & --- \\
K25117 & isdB & T2 & Transport & Fe & --- \\
K25118 & isdC & T2 & Transport & Fe & --- \\
K25119 & shp & T2 & Transport & Fe & --- \\
K25130 & fecE & T2 & Transport & Fe & --- \\
K25132 & hmuT & T2 & Transport & Fe & --- \\
K25133 & hmuU & T2 & Transport & Fe & --- \\
K25282 & yclQ & T2 & Transport & Fe & --- \\
K25283 & yclO & T2 & Transport & Fe & --- \\
K25284 & yclN & T2 & Transport & Fe & --- \\
K25285 & yclP & T2 & Transport & Fe & --- \\
K25286 & fagD & T2 & Transport & Fe & --- \\
K25287 & desE & T2 & Transport & Fe & --- \\
K25288 & fhuB & T2 & Transport & Fe & --- \\
K25289 & fhuG & T2 & Transport & Fe & --- \\
K25290 & fhuA & T2 & Transport & Fe & --- \\
K25308 & fhuD & T2 & Transport & Fe & --- \\
K28864 & YSL-PS & T2 & Transport & --- & --- \\
\midrule
K03325 & ACR3 & T1 & Resistance & As, Tl & BacMet2, Fitness \\
K03446 & emrB & T2 & Resistance & Cu, Tl & BacMet2, Fitness \\
K07665 & cusR & T1 & Resistance & Cu, Ni & BacMet2, Fitness \\
K07785 & nrsD & T1 & Resistance & Ni & Fitness \\
K07787 & cusA & T1 & Resistance & Cu, Ag & BacMet2, Fitness \\
K07798 & cusB & T1 & Resistance & Cu, Ag, Zn & BacMet2, Fitness \\
K08365 & merR & T1 & Resistance & Hg, Cu, Zn & BacMet2, Fitness \\
K15725 & czcC & T1 & Resistance & Ni, Zn, Co, Cd & BacMet2, Fitness \\
K15726 & czcA & T1 & Resistance & Co, Ni, Zn, Cd & BacMet2, Fitness \\
K15727 & czcB & T1 & Resistance & Co, Ni, Zn, Cd & BacMet2, Fitness \\
K16264 & czcD & T1 & Resistance & Cd, Co, Zn, Tl & BacMet2, Fitness \\
K17686 & copA & T2 & Resistance & Cu & BacMet2, Fitness \\
K19591 & cueR & T1 & Resistance & Cu, Co, Ni, Zn & BacMet2, Fitness \\
K19594 & gesB & T1 & Resistance & Cu, Au & BacMet2 \\
K19595 & gesA & T1 & Resistance & Cu, Au & BacMet2 \\
\midrule
K00240 & sdhB & T2 & Sensing & Fe, S & KEGG \\
K00245 & frdB & T2 & Sensing & Fe, S & KEGG \\
K02483 & --- & T2 & Sensing & Al, Co, Cu & BacMet2, Fitness \\
K07644 & cusS & T1 & Sensing & Cu & BacMet2, Fitness \\
K08941 & pscB & T2 & Sensing & Fe, S & KEGG \\
K13638 & zntR & T2 & Sensing & Zn, Cu & BacMet2, Fitness \\
K17219 & sreA & T2 & Sensing & Mo, S & KEGG \\
K17220 & sreB & T2 & Sensing & S & KEGG \\
K17221 & sreC & T2 & Sensing & S & KEGG \\
K18367 & CoADR & T2 & Sensing & S & KEGG \\
\midrule
K02225 & cobC1 & T2 & Cofactor & Co & KEGG \\
K03635 & MOCS2B & T2 & Cofactor & Mo & KEGG \\
K03638 & moaB & T2 & Cofactor & Mo & KEGG \\
K03750 & moeA & T2 & Cofactor & Mo & KEGG \\
K03831 & mogA & T2 & Cofactor & Mo & KEGG \\
\midrule
K00380 & cysJ & T2 & Metabolism & S, Co, Cu & KEGG, Fitness \\
\midrule
K02484 & --- & T2 & --- & Al, Cu & BacMet2, Fitness \\
K03673 & dsbA & T2 & --- & Ag, Al, Bi, Zn & BacMet2, Fitness \\
K03981 & dsbC & T2 & --- & Al, Cu & BacMet2, Fitness \\
\end{longtable}

% ============================================================================
% -- Table S3: Double-signal HGT candidate KOs --------------------------------

\begin{table}[H]
\caption{Double-signal HGT candidate KOs: genome-level Fritz \& Purvis D $> 0.2$
\textbf{and} genus-level Pagel's $\lambda < 0.3$.
D was computed on the 18{,}961-tip GTDB genome tree using binary presence/absence;
$\lambda$ was estimated from genus presence fractions on the GTDB genus-level
phylogeny.
Low D indicates phylogenetically dispersed binary presence at the genome scale;
low $\lambda$ indicates low genus-level signal in continuous presence fractions.
KOs are sorted by D (descending).}
\label{tab:double_signal}
\centering
\small
\begin{tabular}{llllrrp{3.5cm}}
\toprule
KO & Gene & Subcategory & Tier & D & $\lambda$ & Metal relevance \\
\midrule
K07785 & \textit{nrsD} & Resistance/Detox. & T2 & 0.821 & 0.089 & Ni resistance (NrsD repressor) \\
K19059 & \textit{merE} & Transport/Homeostasis & T2 & 0.728 & 0.102 & Hg transport (mer operon) \\
K19057 & \textit{merD} & Resistance/Detox. & T2 & 0.701 & 0.165 & Hg resistance (mer operon) \\
K19594 & \textit{gesB} & Resistance/Detox. & T2 & 0.597 & 0.156 & Au/Te resistance \\
K08356 & \textit{aoxB} & Resistance/Detox. & T2 & 0.562 & 0.000 & As(III) oxidase \\
K19595 & \textit{gesA} & Resistance/Detox. & T2 & 0.458 & 0.161 & Au/Te resistance \\
K25119 & \textit{shp} & Transport/Homeostasis & T3 & 0.385 & 0.000 & Fe transport \\
K03897 & \textit{iucD} & Unknown & T3 & 0.354 & 0.237 & Fe siderophore biosynthesis \\
K19592 & \textit{golS} & Sensing/Regulation & T2 & 0.265 & 0.135 & Au/Te sensing \\
K05908 & \textit{doxDA} & Metal-dep. Metabolism & T3 & 0.254 & 0.000 & As oxidation (sox pathway) \\
K08170 & \textit{norB} & Resistance/Detox. & T3 & 0.239 & 0.033 & Cu/NO detoxification \\
K14974 & \textit{nicC} & Transport/Homeostasis & T3 & 0.224 & 0.000 & Ni transport \\
K15585 & \textit{nikB} & Transport/Homeostasis & T3 & 0.202 & 0.000 & Ni transport \\
\bottomrule
\end{tabular}
\end{table}

% -- Table S4: Tl/Al/S metal KO curation rationale ----------------------------

\begin{longtable}{lllllp{4.8cm}}
\caption{Curation rationale for KOs assigned to thallium (Tl), aluminium (Al),
and sulfur (S): elements with fewer dedicated resistance systems whose KOs were
assigned based on documented biochemical function or fitness-screen evidence.
Tier codes: T1 = direct biochemical evidence; T2 = strong biochemical inference;
T2F = fitness-screen evidence; T3 = computational/BacMet annotation.
$n = 153$ KOs (Al $= 34$, Tl $= 45$, S $= 74$).}
\label{tab:tl_al_s}\\
\toprule
KO & Gene & Metal & Subcategory & Tier & Biochemical rationale \\
\midrule
\endfirsthead
\multicolumn{6}{l}{\small\textit{Table S4 continued.}}\\
\toprule
KO & Gene & Metal & Subcategory & Tier & Biochemical rationale \\
\midrule
\endhead
\midrule
\multicolumn{6}{r}{\small\textit{continued on next page}}\\
\endfoot
\bottomrule
\endlastfoot
K02483 & --- & Al & Sensing/Regulation & T2 & Sensing/regulatory response to Al,Co,Cu \\
K02484 & --- & Al & Unknown & T2 & Fitness screen: defect under Al,Cu stress \\
K03673 & \textit{dsbA} & Al & Unknown & T2 & Fitness screen: defect under Ag,Al,Bi,Zn stress \\
K03981 & \textit{dsbC} & Al & Unknown & T2 & Fitness screen: defect under Al,Cu stress \\
K00528 & \textit{fpr} & Al & Unknown & T2F & Fitness screen: defect under Al,Co,Tl,Zn stress \\
K01002 & \textit{mdoB} & Al & Unknown & T2F & Fitness screen: defect under Al,Bi,Fe,Tl stress \\
K01069 & \textit{gloB} & Al & Unknown & T2F & Fitness screen: defect under Al,Tl,Zn stress \\
K01073 & \textit{yigI} & Al & Unknown & T2F & Fitness screen: defect under Al,Co,Cu,Ni,Tl stress \\
K01104 & \textit{E3.1.3.48} & Al & Unknown & T2F & Fitness screen: defect under Al stress \\
K01256 & \textit{pepN} & Al & Unknown & T2F & Fitness screen: defect under Al,Ni stress \\
K01262 & \textit{pepP} & Al & Unknown & T2F & Fitness screen: defect under Al,Cu,Ni stress \\
K01414 & \textit{prlC} & Al & Unknown & T2F & Fitness screen: defect under Al,Cu,Ni stress \\
K01626 & \textit{E2.5.1.54} & Al & Unknown & T2F & Fitness screen: defect under Al stress \\
K01784 & \textit{galE} & Al & Resistance/Detox. & T2F & Resistance/detoxification; fitness defect under Al,Tl,Zn stress \\
K02027 & \textit{ABC.MS.S} & Al & Unknown & T2F & Fitness screen: defect under Al,Cu,Ni stress \\
K02028 & \textit{ABC.PA.A} & Al & Unknown & T2F & Fitness screen: defect under Al,Ni stress \\
K02030 & \textit{ABC.PA.S} & Al & Unknown & T2F & Fitness screen: defect under Al,Ni stress \\
K02424 & \textit{tcyA} & Al & Transport/Homeostasis & T2F & Metal transport/homeostasis (Al,Cu,Ni) \\
K03296 & \textit{TC.HAE1} & Al & Unknown & T2F & Fitness screen: defect under Al,Cu,Ni stress \\
K03316 & \textit{nhaP} & Al & Unknown & T2F & Fitness screen: defect under Al,Tl stress \\
K03760 & \textit{eptA} & Al & Unknown & T2F & Fitness screen: defect under Al stress \\
K03781 & \textit{katE} & Al & Unknown & T2F & Fitness screen: defect under Al,Ni stress \\
K03814 & \textit{mtgA} & Al & Unknown & T2F & Fitness screen: defect under Al,Cr,Mn stress \\
K03885 & \textit{ndh} & Al & Unknown & T2F & Fitness screen: defect under Al,Tl stress \\
K05810 & \textit{LACC1} & Al & Unknown & T2F & Fitness screen: defect under Al,Tl stress \\
K07165 & \textit{fecR} & Al & Unknown & T2F & Fitness screen: defect under Al,Bi,Tl stress \\
K07277 & \textit{SAM50} & Al & Unknown & T2F & Fitness screen: defect under Al,Ni stress \\
K07278 & \textit{tamA} & Al & Unknown & T2F & Fitness screen: defect under Al,Ni stress \\
K07338 & --- & Al & Unknown & T2F & Fitness screen: defect under Al,Co,Ni stress \\
K07636 & \textit{phoR} & Al & Sensing/Regulation & T2F & Sensing/regulatory response to Al,Cu \\
K07666 & \textit{qseB} & Al & Sensing/Regulation & T2F & Sensing/regulatory response to Al,Cu \\
K09800 & \textit{tamB} & Al & Unknown & T2F & Fitness screen: defect under Al,Co,Ni stress \\
K19353 & \textit{eptC} & Al & Unknown & T2F & Fitness screen: defect under Al stress \\
K20534 & \textit{gtrB} & Al & Unknown & T2F & Fitness screen: defect under Al,Tl stress \\
K00240 & \textit{sdhB} & S & Sensing/Regulation & T2 & Sensing/regulatory response to Fe,S \\
K00245 & \textit{frdB} & S & Sensing/Regulation & T2 & Sensing/regulatory response to Fe,S \\
K00380 & \textit{cysJ} & S & Metal-dep. Metabolism & T2 & Metal-dependent metabolic enzyme (S,Co,Cu) \\
K08941 & \textit{pscB} & S & Sensing/Regulation & T2 & Sensing/regulatory response to Fe,S \\
K17219 & \textit{sreA} & S & Sensing/Regulation & T2 & Sensing/regulatory response to Mo,S \\
K17220 & \textit{sreB} & S & Sensing/Regulation & T2 & Sensing/regulatory response to S \\
K17221 & \textit{sreC} & S & Sensing/Regulation & T2 & Sensing/regulatory response to S \\
K18367 & \textit{CoADR} & S & Sensing/Regulation & T2 & Sensing/regulatory response to S \\
K00024 & \textit{mdh} & S & Metal-dep. Metabolism & T3 & Metal-dependent metabolic enzyme (S) \\
K00031 & \textit{IDH1} & S & Metal-dep. Metabolism & T3 & Metal-dependent metabolic enzyme (S) \\
K00169 & \textit{porA} & S & Unknown & T3 & Fitness screen: defect under S stress \\
K00170 & \textit{porB} & S & Unknown & T3 & Fitness screen: defect under S stress \\
K00171 & \textit{porD} & S & Unknown & T3 & Fitness screen: defect under S stress \\
K00172 & \textit{porC} & S & Unknown & T3 & Fitness screen: defect under S stress \\
K00174 & \textit{korA} & S & Unknown & T3 & Fitness screen: defect under S stress \\
K00175 & \textit{korB} & S & Unknown & T3 & Fitness screen: defect under S stress \\
K00176 & \textit{korD} & S & Unknown & T3 & Fitness screen: defect under S stress \\
K00177 & \textit{korC} & S & Unknown & T3 & Fitness screen: defect under S stress \\
K00239 & \textit{sdhA} & S & Metal-dep. Metabolism & T3 & Metal-dependent metabolic enzyme (S) \\
K00241 & \textit{sdhC} & S & Metal-dep. Metabolism & T3 & Metal-dependent metabolic enzyme (S) \\
K00242 & \textit{sdhD} & S & Metal-dep. Metabolism & T3 & Metal-dependent metabolic enzyme (S) \\
K00244 & \textit{frdA} & S & Metal-dep. Metabolism & T3 & Metal-dependent metabolic enzyme (S) \\
K00246 & \textit{frdC} & S & Metal-dep. Metabolism & T3 & Metal-dependent metabolic enzyme (S) \\
K00247 & \textit{frdD} & S & Metal-dep. Metabolism & T3 & Metal-dependent metabolic enzyme (S) \\
K00381 & \textit{cysI} & S & Metal-dep. Metabolism & T3 & Metal-dependent metabolic enzyme (S) \\
K00390 & \textit{cysH} & S & Unknown & T3 & Fitness screen: defect under S stress \\
K00392 & \textit{sir} & S & Metal-dep. Metabolism & T3 & Metal-dependent metabolic enzyme (S) \\
K00860 & \textit{cysC} & S & Unknown & T3 & Fitness screen: defect under S stress \\
K00955 & \textit{cysNC} & S & Unknown & T3 & Fitness screen: defect under S stress \\
K00956 & \textit{cysN} & S & Unknown & T3 & Fitness screen: defect under S stress \\
K00957 & \textit{cysD} & S & Unknown & T3 & Fitness screen: defect under S stress \\
K00958 & \textit{sat} & S & Unknown & T3 & Fitness screen: defect under S stress \\
K01006 & \textit{ppdK} & S & Unknown & T3 & Fitness screen: defect under S stress \\
K01007 & \textit{pps} & S & Unknown & T3 & Fitness screen: defect under S stress \\
K01595 & \textit{ppc} & S & Unknown & T3 & Fitness screen: defect under S stress \\
K01676 & \textit{E4.2.1.2A} & S & Unknown & T3 & Fitness screen: defect under S stress \\
K01677 & \textit{E4.2.1.2AA} & S & Unknown & T3 & Fitness screen: defect under S stress \\
K01678 & \textit{E4.2.1.2AB} & S & Unknown & T3 & Fitness screen: defect under S stress \\
K01679 & \textit{E4.2.1.2B} & S & Unknown & T3 & Fitness screen: defect under S stress \\
K01681 & \textit{ACO} & S & Unknown & T3 & Fitness screen: defect under S stress \\
K01682 & \textit{acnB} & S & Unknown & T3 & Fitness screen: defect under S stress \\
K01902 & \textit{sucD} & S & Unknown & T3 & Fitness screen: defect under S stress \\
K01903 & \textit{sucC} & S & Unknown & T3 & Fitness screen: defect under S stress \\
K01958 & \textit{PC} & S & Unknown & T3 & Fitness screen: defect under S stress \\
K01959 & \textit{pycA} & S & Unknown & T3 & Fitness screen: defect under S stress \\
K01960 & \textit{pycB} & S & Unknown & T3 & Fitness screen: defect under S stress \\
K02045 & \textit{cysA} & S & Unknown & T3 & Fitness screen: defect under S stress \\
K02046 & \textit{cysU} & S & Transport/Homeostasis & T3 & Metal transport/homeostasis (S) \\
K02047 & \textit{cysW} & S & Transport/Homeostasis & T3 & Metal transport/homeostasis (S) \\
K02048 & \textit{cysP} & S & Unknown & T3 & Fitness screen: defect under S stress \\
K03737 & \textit{por} & S & Unknown & T3 & Fitness screen: defect under S stress \\
K05908 & \textit{doxDA} & S & Metal-dep. Metabolism & T3 & Metal-dependent metabolic enzyme (S) \\
K08940 & \textit{pscA} & S & Unknown & T3 & Fitness screen: defect under S stress \\
K08942 & \textit{pscC} & S & Unknown & T3 & Fitness screen: defect under S stress \\
K08943 & \textit{pscD} & S & Unknown & T3 & Fitness screen: defect under S stress \\
K13811 & \textit{PAPSS} & S & Unknown & T3 & Fitness screen: defect under S stress \\
K15230 & \textit{aclA} & S & Unknown & T3 & Fitness screen: defect under S stress \\
K15231 & \textit{aclB} & S & Unknown & T3 & Fitness screen: defect under S stress \\
K15232 & \textit{ccsA} & S & Unknown & T3 & Fitness screen: defect under S stress \\
K15233 & \textit{ccsB} & S & Unknown & T3 & Fitness screen: defect under S stress \\
K15234 & \textit{ccl} & S & Unknown & T3 & Fitness screen: defect under S stress \\
K16936 & \textit{doxA} & S & Metal-dep. Metabolism & T3 & Metal-dependent metabolic enzyme (S) \\
K16937 & \textit{doxD} & S & Metal-dep. Metabolism & T3 & Metal-dependent metabolic enzyme (S) \\
K17222 & \textit{soxA} & S & Unknown & T3 & Fitness screen: defect under S stress \\
K17223 & \textit{soxX} & S & Unknown & T3 & Fitness screen: defect under S stress \\
K17224 & \textit{soxB} & S & Unknown & T3 & Fitness screen: defect under S stress \\
K17225 & \textit{soxC} & S & Metal-dep. Metabolism & T3 & Metal-dependent metabolic enzyme (S) \\
K17226 & \textit{soxY} & S & Sensing/Regulation & T3 & Sensing/regulatory response to S \\
K17227 & \textit{soxZ} & S & Sensing/Regulation & T3 & Sensing/regulatory response to S \\
K17229 & \textit{fccB} & S & Metal-dep. Metabolism & T3 & Metal-dependent metabolic enzyme (S) \\
K17230 & \textit{fccA} & S & Metal-dep. Metabolism & T3 & Metal-dependent metabolic enzyme (S) \\
K19713 & \textit{tsdA} & S & Metal-dep. Metabolism & T3 & Metal-dependent metabolic enzyme (S) \\
K22622 & \textit{soxD} & S & Unknown & T3 & Fitness screen: defect under S stress \\
K27925 & \textit{tetH} & S & Unknown & T3 & Fitness screen: defect under S stress \\
K03325 & \textit{ACR3} & Tl & Resistance/Detox. & T1 & Resistance/detoxification; fitness defect under As,Tl stress \\
K16264 & \textit{czcD} & Tl & Resistance/Detox. & T1 & Resistance/detoxification; fitness defect under Cd,Co,Zn,Tl stress \\
K03446 & \textit{emrB} & Tl & Resistance/Detox. & T2 & Resistance/detoxification; fitness defect under Cu,Tl stress \\
K00059 & \textit{fabG} & Tl & Unknown & T2F & Fitness screen: defect under Bi,Ni,Tl stress \\
K00108 & \textit{betA} & Tl & Metal-dep. Metabolism & T2F & Metal-dependent metabolic enzyme (Tl) \\
K00130 & \textit{betB} & Tl & Metal-dep. Metabolism & T2F & Metal-dependent metabolic enzyme (Tl) \\
K00523 & \textit{ascD} & Tl & Unknown & T2F & Fitness screen: defect under Co,Cu,Tl stress \\
K00528 & \textit{fpr} & Tl & Unknown & T2F & Fitness screen: defect under Al,Co,Tl,Zn stress \\
K01002 & \textit{mdoB} & Tl & Unknown & T2F & Fitness screen: defect under Al,Bi,Fe,Tl stress \\
K01069 & \textit{gloB} & Tl & Unknown & T2F & Fitness screen: defect under Al,Tl,Zn stress \\
K01073 & \textit{yigI} & Tl & Unknown & T2F & Fitness screen: defect under Al,Co,Cu,Ni,Tl stress \\
K01278 & \textit{DPP4} & Tl & Unknown & T2F & Fitness screen: defect under Cu,Tl stress \\
K01782 & \textit{fadJ} & Tl & Resistance/Detox. & T2F & Resistance/detoxification; fitness defect under Co,Tl stress \\
K01784 & \textit{galE} & Tl & Resistance/Detox. & T2F & Resistance/detoxification; fitness defect under Al,Tl,Zn stress \\
K01814 & \textit{hisA} & Tl & Resistance/Detox. & T2F & Resistance/detoxification; fitness defect under Co,Tl stress \\
K01990 & \textit{ABC-2.A} & Tl & Unknown & T2F & Fitness screen: defect under Co,Fe,Ni,Tl stress \\
K02072 & \textit{metI} & Tl & Transport/Homeostasis & T2F & Metal transport/homeostasis (Cr,Cu,Tl) \\
K02167 & \textit{betI} & Tl & Sensing/Regulation & T2F & Sensing/regulatory response to Tl \\
K02493 & \textit{hemK} & Tl & Unknown & T2F & Fitness screen: defect under Tl stress \\
K03316 & \textit{nhaP} & Tl & Unknown & T2F & Fitness screen: defect under Al,Tl stress \\
K03328 & \textit{TC.PST} & Tl & Unknown & T2F & Fitness screen: defect under Co,Ni,Tl stress \\
K03439 & \textit{trmB} & Tl & Unknown & T2F & Fitness screen: defect under Tl stress \\
K03442 & \textit{mscS} & Tl & Unknown & T2F & Fitness screen: defect under Co,Tl stress \\
K03453 & \textit{TC.BASS} & Tl & Unknown & T2F & Fitness screen: defect under Tl stress \\
K03568 & \textit{tldD} & Tl & Unknown & T2F & Fitness screen: defect under Ni,Tl stress \\
K03593 & \textit{mrp} & Tl & Unknown & T2F & Fitness screen: defect under Tl stress \\
K03701 & \textit{uvrA} & Tl & Unknown & T2F & Fitness screen: defect under Co,Ni,Tl,Zn stress \\
K03885 & \textit{ndh} & Tl & Unknown & T2F & Fitness screen: defect under Al,Tl stress \\
K05368 & \textit{fre} & Tl & Unknown & T2F & Fitness screen: defect under Cu,Tl stress \\
K05592 & \textit{deaD} & Tl & Unknown & T2F & Fitness screen: defect under Cu,Tl stress \\
K05799 & \textit{pdhR} & Tl & Sensing/Regulation & T2F & Sensing/regulatory response to Tl,Zn \\
K05810 & \textit{LACC1} & Tl & Unknown & T2F & Fitness screen: defect under Al,Tl stress \\
K06179 & \textit{rluC} & Tl & Unknown & T2F & Fitness screen: defect under Ni,Tl stress \\
K06941 & \textit{rlmN} & Tl & Unknown & T2F & Fitness screen: defect under Cu,Tl stress \\
K07058 & --- & Tl & Unknown & T2F & Fitness screen: defect under Co,Ni,Tl stress \\
K07165 & \textit{fecR} & Tl & Unknown & T2F & Fitness screen: defect under Al,Bi,Tl stress \\
K07263 & \textit{pqqL} & Tl & Transport/Homeostasis & T2F & Metal transport/homeostasis (Zn,Co,Tl) \\
K07390 & \textit{grxD} & Tl & Unknown & T2F & Fitness screen: defect under Tl,Zn stress \\
K08300 & \textit{rne} & Tl & Unknown & T2F & Fitness screen: defect under Mn,Tl stress \\
K08301 & \textit{rng} & Tl & Unknown & T2F & Fitness screen: defect under Mn,Tl stress \\
K09771 & \textit{TC.SMR3} & Tl & Resistance/Detox. & T2F & Resistance/detoxification; fitness defect under Tl stress \\
K11645 & \textit{fbaB} & Tl & Unknown & T2F & Fitness screen: defect under Ni,Tl stress \\
K13628 & \textit{iscA} & Tl & Sensing/Regulation & T2F & Sensing/regulatory response to Fe,Tl,Zn \\
K18707 & --- & Tl & Unknown & T2F & Fitness screen: defect under Cu,Tl stress \\
K20534 & \textit{gtrB} & Tl & Unknown & T2F & Fitness screen: defect under Al,Tl stress \\
\end{longtable}

% ============================================================================
\section*{Supplementary Figures}
\addcontentsline{toc}{section}{Supplementary Figures}
\setcounter{figure}{0}
% ============================================================================

% -- Figure S1 ----------------------------------------------------------------
\begin{figure}[p]
\centering
\includegraphics[width=0.95\textwidth]{figures/png/figS01_phylogenetic_tree}
\caption{\textbf{Phylogenetic distribution of genera in the primary dataset.}
Circular phylogenetic cladogram (GTDB r214v1; $n = 2{,}283$ bacterial genera
before pruning; $n = 1{,}574$ with PGLS data, shown with filled tips).
Inner ring: Tier~1+2 per-Mb KO density (red--blue gradient; red = high, blue = low).
Outer ring: standardised Levins' niche breadth $B_\mathrm{std}$ (light = specialist,
dark = generalist).  Arc labels indicate the four focal phyla
(Proteobacteria, Actinobacteria, Firmicutes, Bacteroidetes).}
\label{fig:s1}
\end{figure}

\clearpage
% -- Figure S2 ----------------------------------------------------------------
\begin{figure}[p]
\centering
\includegraphics[width=0.95\textwidth]{figures/png/figS02_genome_size}
\caption{\textbf{Genome-size distribution and its relationship to niche breadth.}
(A) Histogram of mean genome size (Mb) across 1,574 genera in the primary
dataset.  (B)~Scatter plot of mean genome size vs standardised niche breadth
($B_\mathrm{std}$) with PGLS regression line.  (C)~Scatter plot of mean
genome size vs \texttt{ko\_per\_mb\_primary\_z}, demonstrating that per-Mb
normalisation partially but not fully accounts for genome-size effects.
Genome size attenuates the primary $\beta$ by 46.7\%, below the pre-specified
50\% confounding threshold.}
\label{fig:s2}
\end{figure}

\clearpage
% -- Figure S3 ----------------------------------------------------------------
\begin{figure}[p]
\centering
\includegraphics[width=0.95\textwidth]{figures/png/figS03_sensitivity}
\caption{\textbf{Sensitivity analyses: PGLS $\beta$ is consistent across eight
pre-specified checks.}
Forest plot of $\beta$ (with $\pm 1.96$~SE error bars) from the primary analysis
(bacteria, $n = 1{,}574$) and eight sensitivity analyses: OLS ($\lambda = 0$),
Brownian motion ($\lambda = 1$), archaea ($n = 41$, $n_\mathrm{min} = 5$),
genus-count thresholds ($n_\mathrm{min} \geq 50$ and $\geq 150$), raw (unstandardised)
Levins $B$, Australia-only subset, northern hemisphere only, and soil-restricted.
Filled points: $p < 0.05$; open points: $p \geq 0.05$.  Six of eight analyses
are significant and directionally consistent.}
\label{fig:s3}
\end{figure}

\clearpage
% -- Figure S4 ----------------------------------------------------------------
\begin{figure}[p]
\centering
\includegraphics[width=0.95\textwidth]{figures/png/figS04_negative_controls}
\caption{\textbf{Named negative controls relative to the primary metal-gene estimate.}
Forest plot of PGLS $\beta$ for three pre-specified named negative controls
(ribosomal proteins, amino acid biosynthesis, DNA repair) alongside the primary
metal-gene estimate.  All three negative controls are more negative than
the primary estimate, confirming that the housekeeping baseline lies at
$\beta \approx -0.030$ to $-0.035$, and the metal-gene $\beta = -0.021$ is
weaker than the streamlining baseline.  Error bars: $\pm 1.96$~SE.}
\label{fig:s4}
\end{figure}

\clearpage
% -- Figure S5 ----------------------------------------------------------------
\begin{figure}[p]
\centering
\includegraphics[width=0.95\textwidth]{figures/png/figS05_permutation}
\caption{\textbf{Coreness-matched permutation: the primary $\beta$ is not
distinctive relative to matched gene sets.}
Histogram of $\beta$ from 1,000 coreness-matched null KO sets ($n = 140$ KOs each,
matched to the primary set by pan-genome prevalence decile).  Red vertical line
marks the observed $\beta = -0.021$ (empirical $p = 0.298$).  Null distribution:
median $= -0.018$, SD $= 0.007$.  The primary metal-gene $\beta$ falls within the
null distribution, indicating that the streamlining signal is not statistically
unusual for a gene set of this size and prevalence distribution.}
\label{fig:s5}
\end{figure}

\clearpage
% -- Figure S6 ----------------------------------------------------------------
\begin{figure}[p]
\centering
\includegraphics[width=0.95\textwidth]{figures/png/figS06_subcategory}
\caption{\textbf{Functional subcategory comparison across AMR, TCS, and ABC
transporter families.}
Side-by-side forest plots of PGLS $\beta$ for subcategories within four gene
families: metal genes (five functional subcategories within the 140-KO primary
set), AMR (efflux pumps, enzymatic inactivation, target modification), TCS
(sensor kinases, response regulators, phosphotransfer proteins), and ABC
transporters (by substrate class).  Demonstrates that the resistance-null /
cofactor-significant split is specific to the metal-gene set; AMR subcategories
are all positive (anti-streamlining), TCS subcategories are all positive,
while ABC transporter Lipid/LPS subcategory shows the only negative signal
($\beta = -0.030$, $q = 3.3 \times 10^{-5}$).  Error bars: $\pm 1.96$~SE;
filled points: BH-FDR $q < 0.05$.}
\label{fig:s6}
\end{figure}

\clearpage
% -- Figure S7 ----------------------------------------------------------------
\begin{figure}[p]
\centering
\includegraphics[width=0.95\textwidth]{figures/png/figS07_cwm_biome}
\caption{\textbf{Community-weighted metal-gene investment across biomes.}
Bar plot of mean \texttt{ko\_per\_mb\_primary\_z} by MicrobeAtlas
\texttt{Env\_Level\_1} habitat category ($n = 12$ categories with $>100$ samples),
weighted by genus relative abundance within each habitat.  Error bars: $\pm$SE.
Marine habitats show consistently lower community-weighted metal-gene density
than terrestrial or host-associated environments, consistent with selective
pressure from oligotrophic marine conditions favouring extreme genome streamlining.}
\label{fig:s7}
\end{figure}

\clearpage
% -- Figure S8 ----------------------------------------------------------------
\begin{figure}[p]
\centering
\includegraphics[width=0.95\textwidth]{figures/png/figS08_ko_cooccurrence}
\caption{\textbf{Tier co-occurrence heatmap.}
Spearman $\rho$ heatmap of cross-genus co-occurrence between individual
Tier~1 and Tier~2 KO densities at the genus level ($118 \times 118$ KOs
represented; BH-FDR $q < 0.05$ cells highlighted).  KOs are clustered by
functional subcategory (indicated by coloured column/row bars).
Demonstrates the within-panel correlation structure: transport KOs form a
broad positive cluster; resistance KOs and cofactor KOs form distinct modules
weakly correlated with the transport block, justifying their treatment as
distinct analytical subcategories.}
\label{fig:s8}
\end{figure}

\clearpage
% -- Figure S9 ----------------------------------------------------------------
\begin{figure}[p]
\centering
\includegraphics[width=0.95\textwidth]{figures/png/figS09_clade_stratified}
\caption{\textbf{Clade-stratified PGLS: signal within Proteobacteria and
Actinobacteria.}
Forest plot of PGLS $\beta$ for four phylum-stratified analyses
(Proteobacteria, Actinobacteria, Firmicutes, Bacteroidetes) alongside the
primary all-bacteria estimate.  Error bars: $\pm 1.96$~SE; filled points:
BH-FDR $q < 0.05$ across the four phyla.
Cochran's $Q = 3.60$ (df $= 3$, $p = 0.309$; $I^2 = 16.6\%$): no significant
heterogeneity across phyla.  The signal is statistically significant in
Proteobacteria and Actinobacteria; Firmicutes and Bacteroidetes are
directionally consistent but underpowered.}
\label{fig:s9}
\end{figure}

\clearpage
% -- Figure S10 ---------------------------------------------------------------
\begin{figure}[p]
\centering
\includegraphics[width=0.85\textwidth]{figures/png/figS10_heterogeneity}
\caption{\textbf{Heterogeneity test: phylum-level $\beta$ estimates are
statistically consistent.}
Funnel plot of phylum-specific $\beta$ ($x$-axis) vs precision $1/\mathrm{SE}$
($y$-axis) for the four clade-stratified estimates
(Proteobacteria, Actinobacteria, Firmicutes, Bacteroidetes).
Cochran's $Q = 3.60$ (df~$= 3$, $p = 0.309$; $I^2 = 16.6\%$): no significant
heterogeneity; the primary effect is not concentrated in a single clade.}
\label{fig:s10}
\end{figure}

\clearpage
% -- Figure S11 ---------------------------------------------------------------
\begin{figure}[p]
\centering
\includegraphics[width=0.95\textwidth]{figures/png/figS11_mag_quality}
\caption{\textbf{MAG quality metrics and their effect on the primary association.}
(A) Histograms of mean per-genus completeness (\%) and contamination (\%) for
the $n = 1{,}107$ genera with CheckM quality metadata.
(B)~Forest plot of PGLS $\beta$ from four models: primary analysis, primary with
completeness added as covariate, primary with contamination as covariate, and
primary restricted to high-quality MAGs (completeness $\geq 90\%$, contamination
$\leq 5\%$; $n = 511$ genera).  Error bars: $\pm 1.96$~SE.  The primary $\beta$
is robust to MAG quality variation: adding quality covariates leaves $\beta$
unchanged at $-0.023$ ($p \approx 3 \times 10^{-8}$).}
\label{fig:s11}
\end{figure}

\clearpage
% -- Figure S12 ---------------------------------------------------------------
\begin{figure}[p]
\centering
\includegraphics[width=0.95\textwidth]{figures/png/figS12_niche_validation}
\caption{\textbf{Niche-breadth metric validation: $B_\mathrm{std}$ correlates
with EMP and MGnify biome metrics.}
(A)~Scatter plot of MicrobeAtlas $B_\mathrm{std}$ vs EMP-derived Levins'
$B_\mathrm{std}$ for 539 overlapping genera (Spearman $\rho = 0.211$,
$p = 7.7 \times 10^{-7}$).
(B)~Scatter plot of MicrobeAtlas $B_\mathrm{std}$ vs MGnify biome Shannon
diversity per genus ($\rho = +0.063$, $p = 0.046$).
(C)~Bootstrap stability: histogram of $|\Delta B_\mathrm{std}|$ between
the original and mean bootstrap estimate across 100 resamples
($n = 1{,}574$ genera).  The low median $|\Delta|$ confirms that
$B_\mathrm{std}$ is stable at the current MicrobeAtlas sample size.}
\label{fig:s12}
\end{figure}

\clearpage
% -- Figure S13 ---------------------------------------------------------------
\begin{figure}[p]
\centering
\includegraphics[width=0.85\textwidth]{figures/png/figS13_emp_pgls}
\caption{\textbf{EMP PGLS: directionally consistent at $n = 539$.}
Scatter plot of EMP-derived Levins' $B_\mathrm{std}$ vs
\texttt{ko\_per\_mb\_primary\_z} for 539 genera shared between the Earth
Microbiome Project and the primary pangenome dataset.
PGLS regression line (blue): $\beta = -0.019$, $p = 0.099$ (non-significant;
directionally consistent with the primary result).
Primary MicrobeAtlas regression line (grey) overlaid for comparison.
The lower significance reflects the smaller sample size ($n = 539$ vs 1,574)
and the coarser EMPO-2 habitat categorisation used as the niche axis.}
\label{fig:s13}
\end{figure}

\clearpage
% -- Figure S14 ---------------------------------------------------------------
\begin{figure}[p]
\centering
\includegraphics[width=0.95\textwidth]{figures/png/figS14_pfam_qc}
\caption{\textbf{Pfam / InterPro audit of 140 primary KOs.}
Left: pie chart showing classification of 140 KOs by Pfam domain outcome:
metal-binding clan confirmed (10/140, 7.1\%; five InterPro clans: CL0704, CL0344,
CL0486, CL0361, CL0193); metal-processing Pfam but no metal-binding clan
(113/140; e.g., ABC transporter scaffolds PF00005/PF00950, MerR-family sensors);
no Pfam annotation (2/140); no representative bacterial gene in KEGG REST (15/140).
Right: breakdown by functional subcategory, demonstrating that low clan coverage
is expected for transporters (which act via binding-pocket selectivity) but not
for cofactor enzymes (which carry structural metal-binding domains).}
\label{fig:s14}
\end{figure}

\clearpage
% -- Figure S15 ---------------------------------------------------------------
\begin{figure}[p]
\centering
\includegraphics[width=0.95\textwidth]{figures/png/figS15_ausmicrobiome}
\caption{\textbf{AusMicrobiome within-dataset PGLS and composition diagnostics.}
(A)~Scatter plot of $B_\mathrm{std}$ vs \texttt{ko\_per\_mb\_primary\_z} for
482 Australian genera ($\beta = -0.052$, $p = 2.2 \times 10^{-15}$).
(B)~Phylum composition bar chart comparing the global ($n = 1{,}574$) and
Australian ($n = 482$) genus panels; Proteobacteria are enriched and
Firmicutes are depleted in the continental subset.
(C)~Primary analysis restricted to the same 482 genera: $\beta = -0.052$,
confirming that the larger within-dataset $\beta$ reflects a composition effect
(the 482 Australian genera, dominated by Proteobacteria, collectively lie in
a steeper part of the global regression).}
\label{fig:s15}
\end{figure}

\clearpage
% -- Figure S16 ---------------------------------------------------------------
\begin{figure}[p]
\centering
\includegraphics[width=0.95\textwidth]{figures/png/figS16_ngsa_map}
\caption{\textbf{Geochemical replication: spatial distribution of NGSA sampling
sites and genus-level metal concentration assignments.}
Map of Australia showing National Geochemical Survey of Australia (NGSA)
ICP-MS soil sampling sites ($n \approx 1{,}200$ sites) and the 200\,km radius
used for spatial assignment of soil metal concentrations to AusMicrobiome
sampling locations.  Inset: histogram of per-genus mean Cu concentration
(ICP-MS, mg\,kg$^{-1}$) for the 482 genera with NGSA coverage.  Sites are
distributed across all major geological provinces, providing broad coverage of
Australian bedrock diversity.  The 200\,km radius was chosen to maximise
genus coverage while remaining within a geologically coherent region for most
sites.}
\label{fig:s16}
\end{figure}

\end{document}
